\chapter{Assinaturas tipo Schnorr avançadas}
\label{chapter:advanced-schnorr}

Entre as ferramentas do capítulo anterior, construímos uma prova de Schnorr a
partir de uma só relação \(K=kG\). Podemos agora ligar várias relações ao mesmo desafio,
esconder a posição da relação verdadeira num anel e, por fim, tornar as
assinaturas ligáveis. SAG e bLSAG servirão de antepassados conceptuais; MLSAG
explica a evolução histórica de RingCT; CLSAG é a construção usada pelas
transacções activas na versão 16 do protocolo. Esta ordem apresenta a extensão
progressiva do mecanismo sem confundir história com consenso actual.

Em todo o capítulo trabalhamos no subgrupo de ordem prima \(l\) e escrevemos
\(\mathbb Z_l\) para os escalares. As listas passadas a uma função hash têm
uma codificação canónica e não ambígua: incluem a ordem e o número dos seus
elementos. \(\mathcal H_n\) designa uma função hash para escalares e
\(\mathcal H_p\), quando aparecer, uma função hash para pontos. São operações
diferentes.

\section{Igualdade do logaritmo discreto em várias bases}
\label{sec:proofs-discrete-logarithm-multiple-bases}

Suponhamos que Alice publica \(K_1=kJ_1\) e \(K_2=kJ_2\). O primeiro ponto
pode ser uma chave pública \(kG\), enquanto o segundo pode ser o segredo
Diffie--Hellman \(kR\) da secção \ref{DH_exchange_section}. Alice quer provar
duas coisas ao mesmo tempo: conhece o logaritmo discreto e esse logaritmo é o
mesmo nas duas bases. Uma resposta Schnorr comum às duas relações faz
precisamente esse trabalho.

\subsection*{Declaração e testemunha}

Seja \(d\geq 1\). A declaração pública, escrita como listas ordenadas, é
\begin{align*}
\mathcal{J} &= (J_1,\ldots,J_d), &
\mathcal{K} &= (K_1,\ldots,K_d).
\end{align*}
A testemunha privada é um único \(k\in\mathbb Z_l\) tal que
\begin{align*}
K_i=kJ_i\qquad\text{para todo }i\in\{1,\ldots,d\}.
\end{align*}
Cada \(J_i\) e \(K_i\) deve ter uma codificação válida, pertencer ao subgrupo
de ordem \(l\) e não ser a identidade. Sem estas condições, uma base nula
poderia tornar uma das relações vazia de conteúdo.

A transformação Fiat--Shamir inclui ainda um rótulo de domínio e um contexto
\(x\). Numa assinatura, \(x\) contém a mensagem \(M\); numa prova
destinada a outro protocolo, contém a sessão e os objectos a que a prova deve
ficar presa.\footnote{Uma prova deste género torna-se uma assinatura quando a
mensagem \(M\) é incluída no contexto do desafio.}

Assumimos que \(\mathcal H_n\) transforma essa codificação em
\(\mathbb Z_l\). Não se exige que cada entrada tenha uma imagem diferente;
uma função hash tem colisões em princípio. A análise usa-a como oráculo
aleatório e a segurança computacional exige que seja impraticável encontrar
colisões ou prever desafios.\footnote{No código Monero, a operação
correspondente calcula Keccak-256 e reduz os 256 bits módulo \(l\). O escalar
zero pertence legitimamente a \(\mathbb Z_l\); a segurança não depende de
excluir esse resultado, cuja ocorrência é de probabilidade desprezável.}

\subsection*{Prova não interactiva}

\begin{enumerate}
    \item Escolhe-se \(\alpha\in_R\mathbb Z_l\) e calculam-se os compromissos
    \(A_i=\alpha J_i\), para \(i=1,\ldots,d\).
    \item Calcula-se o desafio
    \[
    c=\mathcal H_n(\mathsf{DLEQ},x,
       \mathcal J,\mathcal K,[A_1],\ldots,[A_d]).
    \]
    \item Calcula-se a resposta
    \[
    r=(\alpha-ck)\bmod l.
    \]
    \item Publica-se a prova \((c,r)\). Os compromissos não precisam de ser
    publicados separadamente, porque o verificador os reconstrói.
\end{enumerate}

\subsection*{Verificação}

O verificador rejeita primeiro pontos ou escalares mal codificados, pontos
fora do subgrupo, identidades e dimensões incoerentes. Depois calcula
\begin{align*}
A'_i &= rJ_i+cK_i,\qquad i=1,\ldots,d,\\
c' &= \mathcal H_n(\mathsf{DLEQ},x,
       \mathcal J,\mathcal K,[A'_1],\ldots,[A'_d])
\end{align*}
e aceita apenas se \(c'=c\).

\subsection*{Funciona porque}

Para uma prova honesta, cada reconstrução devolve o compromisso original:
\begin{align*}
A'_i
  &=rJ_i+cK_i\\
  &=(\alpha-ck)J_i+c(kJ_i)\\
  &=\alpha J_i=A_i.
\end{align*}
Logo o verificador recompõe exactamente a mesma entrada de
\(\mathcal H_n\), pelo que
\begin{align*}
\mathcal H_n(\mathsf{DLEQ},x,\mathcal J,\mathcal K,
 [A_1],\ldots,[A_d])
&=\mathcal H_n(\mathsf{DLEQ},x,\mathcal J,\mathcal K,
 [A'_1],\ldots,[A'_d]),\\
c&=c'.
\end{align*}

Com \(d=1\), recuperamos a prova de Schnorr da secção
\ref{sec:schnorr-fiat-shamir}. Para \(d>1\), usar o mesmo \(\alpha\), o mesmo
desafio e a mesma resposta liga as relações. No modelo do oráculo aleatório,
duas transcrições aceitáveis com os mesmos compromissos e desafios distintos
permitem extrair \(k\); esta é a razão formal para falar em prova de
conhecimento. A conclusão depende da dificuldade do logaritmo discreto no
subgrupo e da validação dos pontos, não apenas da igualdade \(c=c'\).

O valor \(\alpha\) nunca pode ser reutilizado com outro desafio. De
\(r=\alpha-ck\) e \(r'=\alpha-c'k\), um observador obtém
\[
k=(r-r')(c'-c)^{-1}\pmod l.
\]
Também por isso, qualquer geração determinística do valor efémero deve ficar
presa a todo o contexto e a todas as chaves públicas.

\section{Uma prova com múltiplas chaves privadas}
\label{sec:multiple_private_keys_in_one_proof}

A construção anterior reutilizava uma testemunha em várias bases. Podemos
fazer o complemento: provar simultaneamente várias relações, cada qual com a
sua testemunha, mas reuni-las num só desafio. É uma composição «E» de
provas de Schnorr: a prova só é válida quando todas as relações o são.

\subsection*{Declaração e testemunhas}

Para \(d\geq1\), sejam
\begin{align*}
\mathcal J&=(J_1,\ldots,J_d),&
\mathcal K&=(K_1,\ldots,K_d).
\end{align*}
A declaração pública afirma que existem testemunhas
\((k_1,\ldots,k_d)\in\mathbb Z_l^d\) tais que
\[
K_i=k_iJ_i\qquad\text{para todo }i\in\{1,\ldots,d\}.
\]
As bases podem repetir-se, e podem mesmo ser todas iguais a
\(G\).\footnote{Bases repetidas seriam redundantes numa prova de igualdade de
logaritmos, mas não aqui: por exemplo, \(K_1=k_1G\) e \(K_2=k_2G\) são duas
relações distintas. A construção não afirma que os \(k_i\) sejam diferentes.}
Tal como antes, o verificador exige listas com a mesma dimensão, codificações
canónicas, pontos não nulos no subgrupo de ordem \(l\) e escalares canónicos.
O contexto público \(x\) inclui a mensagem ou a finalidade da prova.

\subsection*{Prova não interactiva}

\begin{enumerate}
    \item Escolhem-se valores efémeros independentes
    \(\alpha_i\in_R\mathbb Z_l\) e calculam-se
    \[
    A_i=\alpha_iJ_i,\qquad i=1,\ldots,d.
    \]
    \item Calcula-se o único desafio
    \[
    c=\mathcal H_n(\mathsf{Schnorr\text{-}AND},x,
      \mathcal J,\mathcal K,[A_1],\ldots,[A_d]).
    \]
    \item Calcula-se uma resposta por testemunha:
    \[
    r_i=(\alpha_i-ck_i)\bmod l,\qquad i=1,\ldots,d.
    \]
    \item Publica-se
    \[
    (c,r_1,\ldots,r_d).
    \]
\end{enumerate}

\subsection*{Verificação}

Depois de validar a declaração e a prova, o verificador reconstrói
\begin{align*}
A'_i&=r_iJ_i+cK_i,\qquad i=1,\ldots,d,\\
c'&=\mathcal H_n(\mathsf{Schnorr\text{-}AND},x,
  \mathcal J,\mathcal K,[A'_1],\ldots,[A'_d])
\end{align*}
e aceita apenas se \(c'=c\).

\subsection*{Funciona porque}

Para cada índice,
\begin{align*}
r_iJ_i+cK_i
  &=(\alpha_i-ck_i)J_i+c(k_iJ_i)\\
  &=\alpha_iJ_i=A_i.
\end{align*}
Assim, todos os argumentos reconstruídos são iguais aos argumentos usados
pelo provador. Abreviando por
\(\mathcal T=(\mathsf{Schnorr\text{-}AND},x,\mathcal J,\mathcal K)\)
o prefixo comum da transcrição, obtemos
\begin{align*}
c'
 &=\mathcal H_n(\mathcal T,[A'_1],\ldots,[A'_d])\\
 &=\mathcal H_n(\mathcal T,[A_1],\ldots,[A_d])=c.
\end{align*}

O desafio partilhado poupa \(d-1\) escalares em comparação com \(d\) provas
Fiat--Shamir independentes; não junta as testemunhas numa só. No modelo do
oráculo aleatório, duas transcrições com os mesmos compromissos e desafios
distintos permitem extrair cada
\[
k_i=(r_i-r'_i)(c'-c)^{-1}\pmod l.
\]
Por conseguinte, a alegação de conhecimento de todas as chaves depende da
dificuldade do logaritmo discreto e de a transcrição ligar o domínio, o
contexto, todas as bases, todas as chaves e todos os compromissos.

Cada \(\alpha_i\) deve ser imprevisível e não pode reaparecer sob outro
desafio: reutilizar apenas um deles basta para revelar o respectivo \(k_i\).
Índices trocados, listas truncadas ou uma codificação ambígua mudariam a
declaração; por isso, dimensões e ordem são parte do que se autentica.

\section{Assinaturas espontâneas anónimas de grupo}
\label{SAG_section}

As assinaturas de grupo de Chaum e van Heyst recorriam a uma entidade de
confiança para formar o grupo e, em certas variantes, para resolver
disputas \cite{Chaum:1991:GS:1754868.1754897}. Uma assinatura em anel segue
outra filosofia: o signatário escolhe espontaneamente um conjunto de chaves
públicas e não precisa da autorização nem da colaboração dos respectivos
donos. A construção foi introduzida por Rivest, Shamir e Tauman
\cite{rivest-leak-secret}; Liu, Wei e Wong desenvolveram depois a variante
ligável LSAG \cite{Liu2004}. Começaremos por SAG, sem ligação, porque o seu
anel de desafios será reutilizado nas construções seguintes.

\subsection*{Declaração, testemunha e índice}

Seja \(M\) a mensagem e seja
\[
\mathcal R=(K_1,\ldots,K_n),\qquad n\geq1,
\]
um anel ordenado de chaves públicas distintas. O signatário conhece um índice
secreto \(\pi\in\{1,\ldots,n\}\) e a testemunha
\(k_\pi\in\mathbb Z_l\) que satisfaz
\[
K_\pi=k_\pi G.
\]
A declaração é uma disjunção: «conheço a chave privada de \emph{algum}
membro de \(\mathcal R\)». A posição \(\pi\) não faz parte da assinatura.
Todas as chaves devem ter codificações canónicas, ser pontos não nulos do
subgrupo de ordem \(l\), e a ordem e o comprimento do anel ficam presos à
transcrição.

Para simplificar os índices, convencionamos que \(c_{n+1}=c_1\). A função
\(\mathcal H_n\) recebe um rótulo de domínio, o anel inteiro, a mensagem e o
compromisso da ronda. O chamado \emph{prefixo de chave} é precisamente esta
inclusão explícita de \(\mathcal R\) no desafio; impede que a mesma
transcrição seja destacada de um anel e apresentada como pertencendo a outro.

\subsection*{Assinatura}

\begin{enumerate}
    \item Escolhe-se \(\alpha\in_R\mathbb Z_l\). Para cada
    \(i\neq\pi\), escolhe-se ainda uma resposta simulada
    \(r_i\in_R\mathbb Z_l\).

    \item Inicia-se a ronda posterior ao signatário:
    \[
    c_{\pi+1}
      =\mathcal H_n(\mathsf{SAG},\mathcal R,M,[\alpha G]).
    \]

    \item Percorrem-se, por ordem cíclica,
    \(i=\pi+1,\pi+2,\ldots,n,1,\ldots,\pi-1\). Em cada posição calcula-se
    \begin{align*}
    L_i&=r_iG+c_iK_i,\\
    c_{i+1}
      &=\mathcal H_n(\mathsf{SAG},\mathcal R,M,[L_i]).
    \end{align*}

    \item Fecha-se o anel com a única resposta que não foi simulada:
    \[
    r_\pi=(\alpha-c_\pi k_\pi)\bmod l.
    \]
\end{enumerate}

A assinatura é
\[
\sigma=(c_1,r_1,\ldots,r_n).
\]
O anel e a mensagem têm de acompanhar a verificação, mas podem estar noutro
campo do protocolo; não é necessário repeti-los dentro desta lista de
escalares.

\subsection*{Verificação}

O verificador confirma \(n\geq1\), o número exacto de respostas e a
canonicidade de \(c_1,r_1,\ldots,r_n\), além de validar todas as chaves do
anel. Depois põe \(\widehat c_1=c_1\) e, para
\(i=1,\ldots,n\), calcula
\begin{align*}
\widehat L_i&=r_iG+\widehat c_iK_i,\\
\widehat c_{i+1}
  &=\mathcal H_n(\mathsf{SAG},\mathcal R,M,[\widehat L_i]).
\end{align*}
No último passo, \(\widehat c_{n+1}\) é o desafio reconstruído para a primeira
posição. Aceita-se se e só se
\[
\widehat c_{n+1}=c_1.
\]
Comparar o desafio final com o inicial, e não com o desafio da ronda
anterior, é o que testa o fecho circular.

\subsection*{Funciona porque}

Nas posições simuladas, o verificador repete literalmente os cálculos do
signatário. Na posição verdadeira,
\begin{align*}
\widehat L_\pi
  &=r_\pi G+c_\pi K_\pi\\
  &=(\alpha-c_\pi k_\pi)G+c_\pi(k_\pi G)\\
  &=\alpha G.
\end{align*}
Assim, também essa ronda devolve o compromisso inicial e a cadeia inteira
fecha.

Vejamos o antigo exemplo de três membros, agora com os índices bem à vista.
Se \(\pi=2\), o signatário escolhe \(\alpha,r_1,r_3\) e calcula
\begin{align*}
c_3&=\mathcal H_n(\mathsf{SAG},\mathcal R,M,[\alpha G]),\\
c_1&=\mathcal H_n(\mathsf{SAG},\mathcal R,M,
                  [r_3G+c_3K_3]),\\
c_2&=\mathcal H_n(\mathsf{SAG},\mathcal R,M,
                  [r_1G+c_1K_1]),\\
r_2&=(\alpha-c_2k_2)\bmod l.
\end{align*}
Como
\[
r_2G+c_2K_2=(\alpha-c_2k_2)G+c_2(k_2G)=\alpha G,
\]
o verificador obtém sucessivamente
\begin{align*}
\widehat c_2&=\mathcal H_n(\mathsf{SAG},\mathcal R,M,
                           [r_1G+c_1K_1]),\\
\widehat c_3&=\mathcal H_n(\mathsf{SAG},\mathcal R,M,
                           [r_2G+\widehat c_2K_2]),\\
\widehat c_1&=\mathcal H_n(\mathsf{SAG},\mathcal R,M,
                           [r_3G+\widehat c_3K_3])=c_1.
\end{align*}

\subsection*{O que a assinatura garante}

Sob a dificuldade do logaritmo discreto e no modelo do oráculo aleatório, a
cadeia prova conhecimento de uma das chaves privadas sem revelar qual. A
distribuição das respostas simuladas deve ser indistinguível da resposta
verdadeira; valores efémeros enviesados, previsíveis ou reutilizados podem
denunciar \(\pi\) e, no caso de reutilização, revelar \(k_\pi\). A
infalsificabilidade exige ainda que mensagem, domínio, anel completo e
compromissos sejam ligados ao hash e que chaves malformadas ou de pequena
ordem sejam rejeitadas.

O valor \(n\) é apenas um limite superior para a cardinalidade do conjunto de
anonimidade que a criptografia proporciona. Se os membros forem escolhidos de
maneira reconhecível ou se informação externa excluir alguns deles, a
ambiguidade efectiva será menor. Para \(n=1\), SAG
reduz-se essencialmente a Schnorr e não oferece anonimidade. Finalmente, SAG
não é ligável: duas assinaturas válidas não revelam, por si só, se usaram a
mesma chave. A imagem de chave da próxima secção acrescentará essa
propriedade.

\section{Assinaturas em anel ligáveis de Adam Back}
\label{blsag_note}

SAG esconde qual membro assinou, mas permite repetir a mesma chave sem deixar
um marcador público. Uma assinatura em anel \emph{ligável} acrescenta esse
marcador: a imagem de chave. O objectivo não é revelar o signatário; é permitir
que qualquer pessoa reconheça uma segunda utilização da mesma chave de uso
único.

LSAG \cite{Liu2004} ligava assinaturas relativas ao mesmo anel ordenado. A
variante que estudaremos, habitualmente chamada bLSAG, torna a imagem
independente dos chamarizes escolhidos. A adaptação foi exposta em
\cite{MRL-0005-ringct} a partir de uma observação de Adam Back
\cite{AdamBack-ring-efficiency} sobre a assinatura CryptoNote
\cite{cryptoNoteWhitePaper}, por sua vez relacionada com o trabalho de
Fujisaki e Suzuki \cite{Fujisaki2007}. É um antepassado conceptual: não é a
assinatura usada pelas transacções Monero activas.

\subsection*{Propriedades pretendidas}

\begin{description}
    \item[Ambiguidade do signatário] Uma assinatura válida demonstra que
    alguém conhece a chave privada de um membro do anel, sem identificar a
    posição, salvo com probabilidade desprezável.
    \footnote{\label{anonymity_note}O anel dá apenas um limite superior para o
    conjunto de anonimidade. Chamarizes escolhidos de forma reconhecível,
    informação temporal ou análise de outras transacções podem excluir
    membros; aumentar \(n\) não transforma automaticamente todos os membros
    em candidatos igualmente plausíveis.}

    \item[Ligação] Duas assinaturas válidas feitas com a mesma chave pública
    de uso único produzem a mesma imagem de chave, mesmo que as mensagens e os
    restantes membros dos anéis sejam diferentes. Reutilizar uma chave apenas
    como chamariz não cria uma ligação.

    \item[Infalsificabilidade] Sem conhecer a chave privada de algum membro,
    um adversário não deve conseguir produzir uma assinatura válida, mesmo
    depois de escolher mensagens e anéis e observar assinaturas legítimas.
\end{description}

Estas propriedades são afirmações criptográficas no modelo do oráculo
aleatório, sob as hipóteses indicadas adiante. Uma imagem repetida não prova,
sozinha, propriedade civil, intenção nem validade de um gasto: é preciso
validar a assinatura e todas as restantes regras do protocolo.

\subsection*{Hash para ponto e imagem de chave}

Além de \(\mathcal H_n\), que devolve escalares, precisamos de
\[
\mathcal H_p:\{0,1\}^*\longrightarrow\mathbb G,
\]
onde \(\mathbb G\) é o subgrupo de ordem prima \(l\). O resultado deve ser um
ponto não nulo e a sua relação de logaritmo discreto com \(G\) não deve ser
conhecida. Não bastaria calcular \(\mathcal H_n(K)G\): como esse multiplicador
é público, qualquer observador obteria
\(\mathcal H_n(K)K\) sem conhecer a chave privada.
\footnote{A família Monero historicamente mapeia a codificação da chave para
um ponto e limpa o cofactor; veja-se a descrição em
\cite{hashtopoint-writeup} e a origem citada em
\cite{hashtopoint-original-paper}. Hash para ponto e hash para escalar não
são operações permutáveis.}

Para o membro verdadeiro \(K_\pi=k_\pi G\), define-se
\[
I=k_\pi\mathcal H_p(K_\pi).
\]
\(I\) é a imagem de chave. Sendo determinística para a relação
\((K_\pi,k_\pi)\), reaparece quando essa chave de uso único volta a assinar;
como a base depende de \(K_\pi\), chaves de uso único diferentes da mesma
carteira não ficam linearmente ligadas.

\subsection*{Declaração, testemunha e assinatura}

Sejam a mensagem \(M\), o anel ordenado
\(\mathcal R=(K_1,\ldots,K_n)\), com \(n\geq1\), o índice secreto
\(\pi\in\{1,\ldots,n\}\) e a testemunha \(k_\pi\) tais que
\(K_\pi=k_\pi G\). Convencionamos novamente \(c_{n+1}=c_1\).

\begin{enumerate}
    \item Calcula-se a imagem
    \(I=k_\pi\mathcal H_p(K_\pi)\).

    \item Escolhem-se \(\alpha\in_R\mathbb Z_l\) e
    \(r_i\in_R\mathbb Z_l\) para todo \(i\neq\pi\).

    \item Inicia-se a ronda seguinte ao signatário:
    \[
    c_{\pi+1}=\mathcal H_n(\mathsf{bLSAG},\mathcal R,I,M,
      [\alpha G],[\alpha\mathcal H_p(K_\pi)]).
    \]

    \item Para
    \(i=\pi+1,\ldots,n,1,\ldots,\pi-1\), calculam-se
    \begin{align*}
    L_i&=r_iG+c_iK_i,\\
    R_i&=r_i\mathcal H_p(K_i)+c_iI,\\
    c_{i+1}
      &=\mathcal H_n(\mathsf{bLSAG},\mathcal R,I,M,
                     [L_i],[R_i]).
    \end{align*}

    \item Fecha-se o anel com
    \[
    r_\pi=(\alpha-c_\pi k_\pi)\bmod l.
    \]
\end{enumerate}

A assinatura é
\[
\sigma=(I,c_1,r_1,\ldots,r_n).
\]
Incluímos \(\mathcal R\) e \(I\) no prefixo de cada desafio. A apresentação
original de bLSAG omitia o anel nesse hash; o prefixo explícito é a práctica
prudente para assinaturas deste género e evita depender de uma interpretação
implícita da transcrição \cite{key-prefix-paper}.

\subsection*{Verificação}

O verificador confirma as dimensões, valida todos os escalares e chaves
públicas e exige
\[
I\neq0\qquad\text{e}\qquad lI=0.
\]
Depois põe \(\widehat c_1=c_1\) e, para \(i=1,\ldots,n\), calcula
\begin{align*}
\widehat L_i&=r_iG+\widehat c_iK_i,\\
\widehat R_i&=r_i\mathcal H_p(K_i)+\widehat c_iI,\\
\widehat c_{i+1}
  &=\mathcal H_n(\mathsf{bLSAG},\mathcal R,I,M,
                 [\widehat L_i],[\widehat R_i]).
\end{align*}
Aceita apenas se
\[
\widehat c_{n+1}=c_1.
\]

\marginnote{\raggedright\tiny
\nolinkurl{src/cryptonote_core/cryptonote_core.cpp}\\
\nolinkurl{check_tx_inputs_keyimages_domain()}}
A condição \(lI=0\), acompanhada da rejeição da identidade, prova que a
imagem pertence ao subgrupo de ordem \(l\). Isto não é o mesmo que
multiplicar simplesmente a equação de verificação pelo cofactor. Se se
aceitasse \(I+T\), com \(T\) um ponto de pequena ordem \(t\), um atacante
poderia explorar a componente de torção e procurar desafios com resíduos
convenientes módulo \(t\), criando variantes que escapassem à ligação. Para
um desafio divisível por \(t\), por exemplo,
\begin{align*}
c(I+T)&=cI+cT\\
      &=cI.
\end{align*}
A verificação de subgrupo elimina todas essas variantes. Esta falha existiu
nos primeiros anos do Monero e foi corrigida antes de ser explorada, nas
regras da versão 5 \cite{key-image-bug}.

\subsection*{Funciona porque}

Na posição verdadeira, as duas reconstruções devolvem os compromissos de
\(\alpha\):
\begin{align*}
\widehat L_\pi
  &=r_\pi G+c_\pi K_\pi\\
  &=(\alpha-c_\pi k_\pi)G+c_\pi(k_\pi G)
   =\alpha G,\\
\widehat R_\pi
  &=r_\pi\mathcal H_p(K_\pi)+c_\pi I\\
  &=(\alpha-c_\pi k_\pi)\mathcal H_p(K_\pi)
    +c_\pi k_\pi\mathcal H_p(K_\pi)\\
  &=\alpha\mathcal H_p(K_\pi).
\end{align*}
As outras rondas foram simuladas exactamente como serão verificadas; portanto
a cadeia fecha no \(c_1\) publicado.

\subsection*{Ligação e hipóteses de segurança}

Dadas duas assinaturas válidas
\begin{align*}
\sigma&=(I,c_1,r_1,\ldots,r_n),\\
\sigma'&=(I',c'_1,r'_1,\ldots,r'_{n'}),
\end{align*}
o teste de ligação é simplesmente \(I=I'\). Para uma mesma chave de uso único,
ambas calculam
\[
I=k_\pi\mathcal H_p(K_\pi),
\]
logo a igualdade é inevitável. No sentido inverso, concluir que imagens
iguais correspondem à mesma relação depende de a função hash para ponto se
comportar como um oráculo aleatório no subgrupo, de serem impraticáveis
colisões e logaritmos discretos e de todos os pontos terem sido validados.

Há uma hipótese distinta por trás da ambiguidade. Para testar se o candidato
\(K_i=k_iG\) corresponde à imagem \(I=k_i\mathcal H_p(K_i)\), um observador
teria de decidir se
\[
(G,K_i,\mathcal H_p(K_i),I)
\]
é um quádruplo Diffie--Hellman. A dificuldade deste problema decisional
impede o teste. O PCDH pede antes \emph{calcular} \(I\), e o PLD pede recuperar
\(k_i\). Um algoritmo para o PLD resolveria também os dois problemas
Diffie--Hellman, mas a mera dificuldade conhecida do PLD não demonstra, em
geral, a dificuldade decisional ou computacional.

Mesmo quando duas assinaturas estão ligadas, a imagem não revela qual membro
assinou. Se os dois anéis tiverem apenas uma chave em comum, porém, a
intersecção denuncia essa chave; é uma perda de anonimidade causada pelo
contexto dos anéis, não uma quebra da equação. A infalsificabilidade exige,
além do logaritmo discreto, prefixos completos, separação de domínio e
resistência do Fiat--Shamir no modelo do oráculo aleatório. Chaves públicas
malformadas, pontos de torção, respostas não canónicas e comprimentos
incoerentes têm de ser rejeitados antes do fecho da cadeia.

Como em Schnorr e SAG, reutilizar \(\alpha\) sob desafios diferentes revela a
testemunha. Uma fonte de aleatoriedade defeituosa pode também distinguir a
posição verdadeira das respostas simuladas. A assinatura guarda
\(n+1\) escalares e uma imagem de chave, além do anel público.

\section{Assinaturas em anel ligáveis com múltiplas camadas}
\label{sec:MLSAG}

MLSAG generaliza bLSAG a uma matriz de chaves
\cite{MRL-0005-ringct}. Foi uma peça central das primeiras transacções
RingCT, pelo que continua a ser necessária para compreender dados históricos
e a evolução do protocolo. Não é, contudo, a assinatura criada nem admitida
para novas transacções na versão 16: CLSAG foi admitida na versão 13 e tornou-se
obrigatória a partir da versão 14 na revisão Monero aqui consultada.

\subsection*{Matriz, declaração e testemunhas}

Seja \(M\) a mensagem. Fixemos a orientação antes de escrever
índices. Há \(n\geq2\) colunas, uma
por membro do anel, e \(m\geq1\) linhas, uma por relação que o signatário deve
provar. Escrevemos
\[
\mathcal R=(K_{i,j})_
 {\substack{1\leq i\leq n\\1\leq j\leq m}}.
\]
O primeiro índice \(i\) escolhe a \emph{coluna}; o segundo \(j\), a
\emph{linha}. A declaração pública é que existe uma coluna secreta
\(\pi\in\{1,\ldots,n\}\) para a qual o signatário conhece todas as
testemunhas
\[
(k_{\pi,1},\ldots,k_{\pi,m})\in\mathbb Z_l^m,
\qquad K_{\pi,j}=k_{\pi,j}G.
\]
Portanto, o conjunto de anonimidade tem no máximo \(n\) colunas, não
\(nm\) chaves independentes. Todas as relações verdadeiras partilham o mesmo
índice \(\pi\); se uma linha denunciar a coluna, denuncia-as a todas.

Na versão integralmente ligável que apresentamos, cada linha tem uma imagem
\[
I_j=k_{\pi,j}\mathcal H_p(K_{\pi,j}),
\qquad j=1,\ldots,m.
\]
Uma variante permite tornar ligáveis apenas \(q\leq m\) linhas: essas linhas
têm termos \(R_{i,j}\) e imagens \(I_j\); as restantes têm apenas termos
\(L_{i,j}\). A fórmula abaixo corresponde a \(q=m\), que torna visível a
estrutura sem esconder casos num índice adicional.

O verificador exige uma matriz rectangular \(n\times m\), pontos canónicos,
não nulos e no subgrupo de ordem \(l\), imagens também não nulas no mesmo
subgrupo e exactamente \(nm\) respostas. Convencionamos \(c_{n+1}=c_1\) e
abreviamos o prefixo completo por
\[
\Phi=(\mathsf{MLSAG},n,m,\mathcal R,I_1,\ldots,I_m,M).
\]
Esta codificação liga o domínio, as dimensões, cada posição da matriz, as
imagens de chave e a mensagem.

\subsection*{Assinatura}

\begin{enumerate}
    \item Calculam-se as imagens
    \[
    I_j=k_{\pi,j}\mathcal H_p(K_{\pi,j}),
    \qquad j=1,\ldots,m.
    \]

    \item Escolhem-se valores independentes
    \(\alpha_j\in_R\mathbb Z_l\), para todas as linhas, e respostas
    \(r_{i,j}\in_R\mathbb Z_l\), para todas as posições com
    \(i\neq\pi\).

    \item Inicia-se a ronda seguinte à coluna signatária:
    \[
    c_{\pi+1}=\mathcal H_n\bigl(
      \Phi,
      [\alpha_1G],[\alpha_1\mathcal H_p(K_{\pi,1})],
      \ldots,
      [\alpha_mG],[\alpha_m\mathcal H_p(K_{\pi,m})]
    \bigr).
    \]

    \item Para as colunas
    \(i=\pi+1,\ldots,n,1,\ldots,\pi-1\), calculam-se, em cada linha,
    \begin{align*}
    L_{i,j}&=r_{i,j}G+c_iK_{i,j},\\
    R_{i,j}&=r_{i,j}\mathcal H_p(K_{i,j})+c_iI_j,
       \qquad j=1,\ldots,m,
    \end{align*}
    e encadeia-se
    \[
    c_{i+1}=\mathcal H_n\bigl(
      \Phi,[L_{i,1}],[R_{i,1}],\ldots,[L_{i,m}],[R_{i,m}]
    \bigr).
    \]

    \item Na coluna verdadeira, definem-se
    \[
    r_{\pi,j}=(\alpha_j-c_\pi k_{\pi,j})\bmod l,
    \qquad j=1,\ldots,m.
    \]
\end{enumerate}

A assinatura é
\[
\sigma=(I_1,\ldots,I_m,c_1,
 r_{1,1},\ldots,r_{1,m},\ldots,r_{n,1},\ldots,r_{n,m}).
\]

\subsection*{Verificação}

Após as validações de pontos, escalares e dimensões, põe-se
\(\widehat c_1=c_1\). Para cada coluna \(i=1,\ldots,n\) e linha
\(j=1,\ldots,m\), calculam-se
\begin{align*}
\widehat L_{i,j}&=r_{i,j}G+\widehat c_iK_{i,j},\\
\widehat R_{i,j}
 &=r_{i,j}\mathcal H_p(K_{i,j})+\widehat c_iI_j.
\end{align*}
O desafio seguinte é
\[
\widehat c_{i+1}=\mathcal H_n\bigl(
 \Phi,
 [\widehat L_{i,1}],[\widehat R_{i,1}],\ldots,
 [\widehat L_{i,m}],[\widehat R_{i,m}]
\bigr).
\]
Aceita-se apenas se o desafio depois da última coluna fechar a cadeia:
\[
\widehat c_{n+1}=c_1.
\]

\subsection*{Funciona porque}

Nas colunas simuladas, os valores reconstruídos são os que o signatário
usou. Na coluna \(\pi\), para cada \(j\),
\begin{align*}
\widehat L_{\pi,j}
 &=r_{\pi,j}G+c_\pi K_{\pi,j}\\
 &=(\alpha_j-c_\pi k_{\pi,j})G
   +c_\pi(k_{\pi,j}G)
  =\alpha_jG,\\
\widehat R_{\pi,j}
 &=r_{\pi,j}\mathcal H_p(K_{\pi,j})+c_\pi I_j\\
 &=(\alpha_j-c_\pi k_{\pi,j})\mathcal H_p(K_{\pi,j})
   +c_\pi k_{\pi,j}\mathcal H_p(K_{\pi,j})\\
 &=\alpha_j\mathcal H_p(K_{\pi,j}).
\end{align*}
O verificador recompõe, linha a linha e na mesma ordem, o compromisso que
iniciou \(c_{\pi+1}\). Uma só resposta errada basta para alterar esse desafio
e impedir o fecho.

\subsection*{Ligação, anonimidade e infalsificabilidade}

Se alguma testemunha \(k_{\pi,j}\) for reutilizada para a mesma chave pública
de uso único, a imagem \(I_j\) reaparece e liga as assinaturas. Uma chave que
surja apenas como chamariz não produz imagem e não cria ligação. Tal como em
bLSAG, anéis cuja intersecção seja pequena podem transformar uma ligação numa
forte pista sobre \(\pi\); a criptografia não corrige um conjunto de
anonimidade mal formado.

A infalsificabilidade exige conhecimento de \emph{todas} as testemunhas da
coluna escolhida, sob a dificuldade do PLD e no modelo do oráculo aleatório.
A associação de uma imagem \(I_j\) ao candidato \(K_{i,j}\) forma uma
instância decisional de Diffie--Hellman
\((G,K_{i,j},\mathcal H_p(K_{i,j}),I_j)\); por isso, a ambiguidade requer que
esse teste decisional seja difícil. Esta hipótese não deve ser confundida com
o PCDH, que pede calcular o quarto ponto, nem com o PLD, que pede recuperar o
escalar. Em grupos gerais, uma destas dificuldades não prova automaticamente
as outras.

Imagens e chaves fora do subgrupo, identidades, matrizes não rectangulares,
ordens de linhas divergentes, desafios ou respostas não canónicos e prefixos
incompletos invalidam o argumento de segurança. Reutilizar
\(\alpha_j\) sob desafios diferentes revela \(k_{\pi,j}\), mesmo que os
outros valores efémeros sejam novos.

\subsection*{Requisitos de espaço}

Na variante com todas as linhas ligáveis, guardam-se \(mn+1\) escalares e
\(m\) imagens de chave; a matriz contém \(mn\) chaves públicas. A assinatura
cresce, portanto, com o produto das linhas pelas colunas. CLSAG conservará um
único ciclo de \(n\) respostas ao agregar as relações de cada coluna.

\section{Assinaturas em anel ligáveis concisas}
\label{sec:CLSAG}

CLSAG, de \emph{concise linkable spontaneous anonymous group signature},
conserva a coluna secreta de MLSAG mas agrega as relações dessa coluna antes
de percorrer o anel \cite{MRL-0011-CLSAG}. Em vez de publicar uma resposta
por linha e por coluna, publica uma resposta por coluna. É esta a assinatura
em anel usada pelas transacções activas na versão 16 do Monero.
\footnote{No código consultado, CLSAG é admitida desde a
versão 13 do protocolo e obrigatória para novas transacções RingCT desde a
versão 14. MLSAG continua necessária para interpretar formatos históricos,
mas não é válida para uma nova transacção v16.}

\subsection*{Chaves primárias e auxiliares}

Seja \(M\) a mensagem. Voltamos à matriz
\[
\mathcal R=(K_{i,j})_
 {\substack{1\leq i\leq n\\1\leq j\leq m}},
\]
com \(n\geq1\) colunas e \(m\geq1\) linhas. A linha \(j=1\) contém as chaves
\emph{primárias}; as linhas \(j>1\), as chaves \emph{auxiliares}. O
signatário conhece uma coluna secreta \(\pi\) e todas as testemunhas
\[
K_{\pi,j}=k_{\pi,j}G,\qquad j=1,\ldots,m.
\]
A ligação deve depender apenas da chave primária \(k_{\pi,1}\); as relações
auxiliares têm de ser provadas, mas não devem criar marcadores de gasto
independentes.

Para cada coluna, define-se uma única base de imagem
\[
H_i=\mathcal H_p(K_{i,1}).
\]
As imagens da coluna verdadeira são
\[
I_j=k_{\pi,j}H_\pi,\qquad j=1,\ldots,m.
\]
\(I_1\) é a imagem primária e será usada no teste de ligação; \(I_2,\ldots,I_m\)
são imagens auxiliares necessárias à agregação.

Todos os escalares, pontos, dimensões e codificações são validados como nas
secções anteriores. Em particular, a imagem primária deve ser não nula e
pertencer ao subgrupo de ordem \(l\). O caso \(n=1\) é algebricamente válido,
mas não oferece anonimidade; nas transacções v16 o anel tem exactamente
16 membros.

\subsection*{Pesos de agregação}

Não podemos somar as linhas com coeficientes fixos. Um atacante poderia
escolher uma chave ou imagem auxiliar que cancelasse outra e obter um
agregado legítimo a partir de componentes ilegítimos. Por isso, CLSAG deriva
um peso independente para cada linha:
\[
\mu_j=\mathcal H_n(T_j,m,n,\mathcal R,I_1,\ldots,I_m),
\qquad j=1,\ldots,m.
\]
Os rótulos \(T_j\) separam os domínios e toda a declaração, incluindo as
imagens, fica presa a cada peso \cite{MRL-0011-CLSAG}. Definem-se então
\begin{align*}
W_i&=\sum_{j=1}^{m}\mu_jK_{i,j},\\
\widetilde W&=\sum_{j=1}^{m}\mu_jI_j,\\
w_\pi&=\sum_{j=1}^{m}\mu_jk_{\pi,j}.
\end{align*}
Na coluna verdadeira,
\begin{align*}
W_\pi&=w_\pi G,&
\widetilde W&=w_\pi H_\pi.
\end{align*}
Assim, muitas relações transformam-se numa só relação bLSAG, sem perder a
vinculação às componentes.

\subsection*{Assinatura}

Seja
\[
\Psi=(T_c,m,n,\mathcal R,I_1,\ldots,I_m,M)
\]
o prefixo das rondas, com um domínio \(T_c\) distinto dos domínios de
agregação. Convencionamos \(c_{n+1}=c_1\).

\begin{enumerate}
    \item Calculam-se \(H_i\), as imagens \(I_j\), os pesos \(\mu_j\), as
    chaves agregadas \(W_i\), a imagem agregada \(\widetilde W\) e a
    testemunha agregada \(w_\pi\).

    \item Escolhem-se \(\alpha\in_R\mathbb Z_l\) e
    \(r_i\in_R\mathbb Z_l\) para todo \(i\neq\pi\).

    \item Inicia-se a ronda seguinte à coluna signatária:
    \[
    c_{\pi+1}
      =\mathcal H_n(\Psi,[\alpha G],[\alpha H_\pi]).
    \]

    \item Para
    \(i=\pi+1,\ldots,n,1,\ldots,\pi-1\), calculam-se
    \begin{align*}
    L_i&=r_iG+c_iW_i,\\
    R_i&=r_iH_i+c_i\widetilde W,\\
    c_{i+1}&=\mathcal H_n(\Psi,[L_i],[R_i]).
    \end{align*}

    \item Fecha-se a cadeia com
    \[
    r_\pi=(\alpha-c_\pi w_\pi)\bmod l.
    \]
\end{enumerate}

A assinatura conceptual é
\[
\sigma=(I_1,\ldots,I_m,c_1,r_1,\ldots,r_n).
\]

\subsection*{Verificação}

O verificador recompõe os mesmos pesos, \(W_i\) e \(\widetilde W\). Depois
põe \(\widehat c_1=c_1\) e, para \(i=1,\ldots,n\), calcula
\begin{align*}
\widehat L_i&=r_iG+\widehat c_iW_i,\\
\widehat R_i&=r_iH_i+\widehat c_i\widetilde W,\\
\widehat c_{i+1}
 &=\mathcal H_n(\Psi,[\widehat L_i],[\widehat R_i]).
\end{align*}
Aceita apenas se
\[
\widehat c_{n+1}=c_1.
\]
O desafio que se compara é, uma vez mais, o valor obtido \emph{depois} da
última coluna; trocar \(c_n\), \(c_{n+1}\) e \(c_1\) seria um erro de fronteira
capaz de destruir o teste circular.

\subsection*{Funciona porque}

Na coluna verdadeira, a resposta agregada reconstrói ambos os compromissos:
\begin{align*}
\widehat L_\pi
 &=r_\pi G+c_\pi W_\pi\\
 &=(\alpha-c_\pi w_\pi)G+c_\pi(w_\pi G)
  =\alpha G,\\
\widehat R_\pi
 &=r_\pi H_\pi+c_\pi\widetilde W\\
 &=(\alpha-c_\pi w_\pi)H_\pi+c_\pi(w_\pi H_\pi)
  =\alpha H_\pi.
\end{align*}
As rondas simuladas coincidem por construção; portanto o desafio volta ao
\(c_1\) publicado.

Os pesos são parte da correcção, não apenas uma compressão. Como cada
\(\mu_j\) depende de todas as chaves e imagens, escolher uma componente para
cancelar outra altera os próprios coeficientes. No modelo do oráculo
aleatório, encontrar o ponto fixo necessário para esse ataque tem
probabilidade desprezável.

\subsection*{Especialização usada pelo Monero}

\marginnote{\raggedright\tiny
\nolinkurl{src/ringct/rctSigs.cpp}\\
{\tt CLSAG\_Gen()}\\
{\tt verRctCLSAG\allowbreak Simple()}}
Nas transacções Monero há duas linhas. A primeira contém chaves públicas de
uso único
\[
P_i=p_iG.
\]
A segunda contém diferenças entre o compromisso candidato
\(C_i^{\mathrm{nz}}\) e um compromisso de compensação comum
\(C_{\mathrm{off}}\):
\[
C_i=C_i^{\mathrm{nz}}-C_{\mathrm{off}}=z_iG.
\]
Na coluna verdadeira, o signatário conhece \(p=p_\pi\) e \(z=z_\pi\). Com
\(H_i=\mathcal H_p(P_i)\), calcula
\begin{align*}
I&=pH_\pi,\\
D&=zH_\pi.
\end{align*}
\(I\) é a imagem primária. A assinatura serializa a imagem auxiliar como
\(\overline D=8^{-1}D\); o verificador recupera \(D=8\overline D\). Esta
multiplicação pelo cofactor coloca a imagem auxiliar no subgrupo primo e o
verificador rejeita \(D=0\). A imagem \(I\), que determina a ligação, é
rejeitada se for nula ou se \(lI\neq0\).

Na implementação consultada, os pesos são exactamente
\begin{align*}
\mu_P&=\mathcal H_n(
 \texttt{CLSAG\_agg\_0},
 \boldsymbol P,\boldsymbol C^{\mathrm{nz}},
 I,\overline D,C_{\mathrm{off}}),\\
\mu_C&=\mathcal H_n(
 \texttt{CLSAG\_agg\_1},
 \boldsymbol P,\boldsymbol C^{\mathrm{nz}},
 I,\overline D,C_{\mathrm{off}}).
\end{align*}
As duas relações agregadas são
\begin{align*}
W_i&=\mu_PP_i+\mu_C(C_i^{\mathrm{nz}}-C_{\mathrm{off}}),\\
\widetilde W&=\mu_PI+\mu_CD,
\end{align*}
e a testemunha verdadeira é
\[
w_\pi=\mu_Pp+\mu_Cz.
\]

O hash de cada ronda usa o domínio \texttt{CLSAG\_round} e liga, pela ordem
serializada, todas as \(P_i\), todos os \(C_i^{\mathrm{nz}}\),
\(C_{\mathrm{off}}\), a mensagem da transacção e os dois compromissos
\(L_i,R_i\). \(I\) e \(\overline D\) ficam ligados às rondas através dos pesos
\(\mu_P,\mu_C\). Em notação compacta:
\begin{align*}
L_i&=r_iG+c_i\{\mu_PP_i+\mu_C(C_i^{\mathrm{nz}}-C_{\mathrm{off}})\},\\
R_i&=r_iH_i+c_i(\mu_PI+\mu_CD).
\end{align*}
A assinatura serializada contém \(r_1,\ldots,r_n\), \(c_1\) e
\(\overline D\). \(I\) já aparece como imagem da entrada da transacção e é
reposto na estrutura antes da verificação; não é duplicado no corpo CLSAG.
O verificador da revisão consultada rejeita anéis vazios, números errados de
respostas, \(c_1\) ou respostas não canónicos, \(I=0\), \(8\overline D=0\) e
codificações de pontos inválidas; a validação da transacção exige ainda
\(lI=0\). Cada desafio de ronda reconstruído tem também de ser não nulo.

\subsection*{Ligação e hipóteses de segurança}

O teste de ligação usa somente
\[
I=k_{\pi,1}\mathcal H_p(K_{\pi,1}).
\]
Se a chave primária de uso único assinar de novo, \(I\) reaparece. As imagens
auxiliares não entram nesse teste: existem para provar a consistência da
agregação, e tratá-las como marcadores poderia criar ligações que o protocolo
não pretende.

A infalsificabilidade e o conhecimento da coluna dependem da dificuldade do
PLD, da segurança Fiat--Shamir no modelo do oráculo aleatório e de os pesos
impedirem cancelamentos de chaves. A impossibilidade de calcular uma imagem
a partir de \(G,K,H\) sem a testemunha é uma hipótese do tipo PCDH; esconder a
qual chave \(K_i\) corresponde a \((H_i,I)\) exige a dificuldade do problema
decisional de Diffie--Hellman. Resolver o PLD permitiria resolver os dois
problemas Diffie--Hellman; o inverso, e a equivalência entre as respectivas
dificuldades, não se seguem sem uma redução para o grupo usado.

A unicidade útil de \(I\) exige ainda uma função hash para ponto resistente a
colisões, bases no subgrupo correcto e rejeição de identidade e torção. Uma
imagem repetida só estabelece ligação entre assinaturas válidas; não prova,
isoladamente, que uma saída pertencia a alguém nem que todas as regras de
valor de uma transacção foram cumpridas.

Reutilizar \(\alpha\) em dois desafios revela a testemunha agregada
\(w_\pi\). Se a reutilização ocorrer sob conjuntos de pesos independentes,
várias equações lineares podem ainda revelar as testemunhas individuais.
Respostas previsíveis ou não uniformes podem denunciar \(\pi\). Por isso,
valor efémero, mensagem, matriz, imagens, dimensões, domínios e todos os prefixos
devem ficar ligados à mesma execução.

\subsection*{Requisitos de espaço}

Na forma geral, CLSAG guarda \(n+1\) escalares e \(m\) imagens, usando
\(nm\) chaves públicas. A redução de \(mn+1\) para \(n+1\) escalares é a
vantagem concisa em relação a MLSAG. Na especialização Monero, a imagem
primária acompanha a entrada e a assinatura serializa apenas a imagem
auxiliar, além desses \(n+1\) escalares.
